\documentclass[11pt,a4paper]{article}

\usepackage[utf8]{inputenc}
\usepackage[T1]{fontenc}
\usepackage[spanish,es-noquoting]{babel}
\usepackage[dvipsnames,table]{xcolor}
\usepackage{amsmath,amssymb}
\usepackage{booktabs}
\usepackage{multirow}
\usepackage{diagbox}
\usepackage{geometry}
\usepackage{caption}
\usepackage{natbib}
\usepackage{graphicx}
\usepackage{hyperref}

\geometry{margin=2.5cm}
\bibliographystyle{apalike}
\hypersetup{colorlinks=true, linkcolor=NavyBlue, citecolor=ForestGreen, urlcolor=NavyBlue}

\definecolor{COLESP}{HTML}{2A78D6}
\definecolor{COLARG}{HTML}{E34948}
\definecolor{COLNEU}{HTML}{888780}

\newcommand{\xG}{\mathrm{xG}}
\newcommand{\E}{\mathbb{E}}

\title{\textbf{Un modelo de \textit{expected goals} con ajuste por rival\\ para la final del Mundial 2026: España vs.\ Argentina}\\[0.4em]
\large Estimación de fuerzas de ataque y defensa por máxima verosimilitud\\ sobre datos de xG a nivel de disparo}
\author{Nestor Lopez}
\date{18 de julio de 2026}

\begin{document}
\maketitle

\begin{abstract}
\noindent
Se construye un modelo probabilístico para la final de la Copa Mundial de la FIFA 2026 (España vs.\ Argentina, MetLife Stadium, 19 de julio de 2026) a partir de los 102 partidos ya disputados del torneo. La estimación se apoya en datos de \textit{expected goals} (xG) a nivel de disparo, obtenidos de Sofascore, que resumen la calidad de cada ocasión con un modelo propietario entrenado sobre grandes volúmenes de datos posicionales. Las fuerzas de generación y de supresión de xG de las 48 selecciones se estiman conjuntamente mediante un modelo log-lineal de Poisson ajustado por la calidad del rival, aprovechando la suficiencia de los totales marginales para prescindir de los datos partido a partido. El modelo revela que el relato dominante ---``el mejor ataque'' (Argentina, 19 goles) contra ``la mejor defensa'' (España, 1 gol encajado)--- describe lo ocurrido pero no lo predice: por xG, ambos finalistas generan peligro de forma casi idéntica ($1{,}92$ de xG por 90 minutos), de modo que la disparidad de goles es varianza de finalización y no diferencia en la creación de ocasiones. La ventaja de España es fundamentalmente defensiva: concede el menor xG por 90 minutos de todo el torneo. El modelo asigna a España una probabilidad de 64{,}9\% de conquistar el título (IC 95\% \textit{bootstrap}: [35{,}6\%; 89{,}1\%]), con $\lambda_{\text{ESP}}=0{,}77$ y $\lambda_{\text{ARG}}=0{,}37$, y anticipa una final muy cerrada (1{,}14 goles esperados frente a una media de torneo de 2{,}73).

\smallskip
\noindent\textbf{Palabras clave:} \textit{expected goals}, modelo Poisson, ajuste por rival, suficiencia, máxima verosimilitud, analítica deportiva.
\end{abstract}

\section{Introducción}

La estimación de probabilidades de resultado en fútbol mediante modelos de Poisson tiene una tradición consolidada desde \citet{maher1982}, refinada por \citet{dixoncoles1997} con una corrección para marcadores bajos y una ponderación temporal. La estructura básica es común: el número de goles de cada equipo se modela como una variable de Poisson cuya media depende multiplicativamente de un parámetro de ataque del anotador, uno de defensa del rival y, cuando corresponde, un término de localía.

La innovación de la última década ha sido reemplazar el gol ---un suceso escaso y ruidoso--- por el \textit{expected goal} (xG) como unidad de análisis. El xG asigna a cada disparo la probabilidad de que termine en gol, estimada a partir de la geometría del remate, la parte del cuerpo, el tipo de jugada y, en los modelos más avanzados, la posición de los defensores \citep{anzerbauer2021}. Su ventaja sobre el recuento de goles es la estabilidad: mide la calidad del proceso ofensivo con menos varianza que el resultado final \citep{robberechts2020}. Esto lo hace especialmente adecuado para muestras cortas, donde las tasas de gol crudas resultan poco fiables.

El presente informe aplica esta idea a un problema con tres particularidades que condicionan las decisiones metodológicas:

\begin{enumerate}
\item \textbf{Muestra muy corta.} Cada finalista ha disputado 7 partidos. Estimar tasas de gol por máxima verosimilitud sobre 7 observaciones produce estimadores de varianza alta y, en casos extremos, degenerados.
\item \textbf{Sede neutral.} El MetLife Stadium no es campo de ninguno de los dos finalistas, de modo que el término de localía se anula por construcción.
\item \textbf{Camino desigual.} España y Argentina no enfrentaron rivales de dificultad equivalente, lo que exige corregir las tasas observadas por la calidad de los adversarios.
\end{enumerate}

El objetivo no es únicamente producir un número, sino documentar por qué el relato periodístico dominante ---``mejor ataque contra mejor defensa''--- tiene escaso valor predictivo, y qué estructura estadística permite corregirlo.

\subsection{Contexto del partido}

España venció a Francia por 2--0 el 14 de julio; Argentina venció a Inglaterra por 2--1 el 15 de julio. La final se disputa el domingo 19 de julio de 2026 en el MetLife Stadium (East Rutherford, Nueva Jersey). Es la primera vez que los campeones vigentes de Europa y de América disputan la final del Mundial. En el historial, ambas selecciones se han enfrentado 14 veces con seis victorias por lado y dos empates: el registro histórico no aporta información discriminante.

\section{Datos}

\subsection{Fuente y extracción}

Los datos de xG proceden de \textbf{Sofascore}, que publica el mapa de disparos (\textit{shotmap}) de cada partido con un valor de xG y de xGOT (\textit{expected goals on target}) por disparo, junto con coordenadas, parte del cuerpo y tipo de jugada. La extracción se realizó mediante el paquete \texttt{ScraperFC} 4.5.0 de Python sobre la competición \texttt{FIFA World Cup}, temporada \texttt{2026}.

La cobertura fue completa: los 102 partidos disputados, con 2.591 disparos y una media de 25{,}4 por encuentro. El recuento de goles ---necesario para separar penales de tiros de juego--- se tomó de FBref, por las razones que se detallan más abajo.

Los goles de penal se tratan por separado. Un modelo de xG describe la fase de juego dinámico; los penales son sucesos de naturaleza distinta, con una probabilidad de conversión estable en torno al 75\% e independiente del rival. Incluirlos en las fuerzas de ataque y defensa contaminaría su interpretación. Se modelan como un término aditivo constante (Sección~\ref{sec:agregacion}).

\subsection{Validación de consistencia}

La suma de xG a favor y en contra sobre las 48 selecciones coincide exactamente (300{,}3 en ambos sentidos), como debe ocurrir cuando cada disparo se asigna a un equipo y a su rival. El cociente entre goles y xG del torneo es $1{,}057$: en agregado se marcó un 5{,}7\% por encima de lo esperado, coherente con la sobre-conversión característica de las fases eliminatorias. La exposición implicada por el calendario (209{,}33 unidades de 90 minutos, con 8 partidos prolongados a la prórroga) reproduce la suma de minutos reportada por la fuente hasta el redondeo.

\subsection{Una precaución de integración}

El campo de resultado de Sofascore incluye los goles de las tandas de penales: un empate a uno resuelto por 4--3 en la tanda figura como ``4--5''. Sumar esos marcadores habría inflado el recuento de goles en 38 unidades respecto a FBref. Se resolvió tomando los goles del recuento de FBref ---limpios de tanda--- y el xG de Sofascore, sin combinar los recuentos de gol de ambas fuentes. El xG, por su parte, no se ve afectado: las tandas de penales no generan disparos de juego y por tanto no contribuyen al xG.

\section{Análisis exploratorio}

\subsection{Goles frente a xG}

El Cuadro~\ref{tab:vsxg} contrasta los goles con el xG para los finalistas. La discrepancia es la clave de todo el análisis.

\begin{table}[htbp]
\centering
\caption{Goles (recuento de FBref, con penales de juego) frente a xG (Sofascore), tras 7 partidos.}
\label{tab:vsxg}
\begin{tabular}{llrrr}
\toprule
& & \textbf{Goles} & \textbf{xG} & \textbf{Diferencia} \\
\midrule
\multirow{2}{*}{Ataque} & Argentina & 18 & 14{,}79 & $+3{,}21$ \\
 & España & 12 & 13{,}47 & $-1{,}47$ \\
\midrule
\multirow{2}{*}{Defensa} & Argentina & 7 & 3{,}94 & $-3{,}06$ \\
 & España & 1 & 2{,}13 & $+1{,}13$ \\
\bottomrule
\end{tabular}
\end{table}

Argentina marcó 18 goles con un xG de 14{,}79: sobrerrindió en ataque por más de tres goles. España marcó 12 con un xG de 13{,}47: finalizó por debajo de lo esperado. En defensa, Argentina encajó 7 goles cuando su xG concedido era de 3{,}94 ---mala suerte o buenas actuaciones rivales---, mientras que España encajó un único gol frente a un xG concedido de 2{,}13. En conjunto, Argentina acumuló unos seis goles de varianza favorable y España finalizó por debajo de su xG ofensivo.

\subsection{La paradoja del gol único}

España encajó un solo gol en todo el torneo. Estimar su tasa defensiva por máxima verosimilitud cruda produce $\hat\lambda_{\text{def}} = 1/7 \approx 0{,}14$ goles por partido. Introducido en un modelo de Poisson, este valor predice que Argentina ---el máximo goleador del torneo--- marcaría aproximadamente $0{,}15$ goles en la final. El resultado es estadísticamente absurdo: no es una señal, es una muestra de 7 partidos con un único evento.

El xG resuelve la paradoja sin necesidad de correcciones \textit{ad hoc}. El xG concedido por España es de $2{,}13$ en 7 partidos, es decir, $0{,}305$ por 90 minutos: bajo, el más bajo del torneo, pero no nulo. El registro de un solo gol combinaba una defensa genuinamente excelente con algo de fortuna. El xG mide la primera y descarta la segunda, ofreciendo un estimador estable de la capacidad defensiva donde el recuento de goles solo ofrecía un valor extremo y engañoso.

\subsection{Calidad de ocasión}

La comparación directa de la calidad ofensiva de los finalistas desmiente el relato dominante. El xG por 90 minutos generado es de $1{,}924$ para España y $1{,}920$ para Argentina: una diferencia de cuatro milésimas. En creación de peligro, ambas selecciones son indistinguibles. Toda la diferencia entre ``19 goles'' y ``12 goles'' proviene de la finalización respecto al valor esperado, no de la generación de ocasiones.

\section{Metodología}
\label{sec:metodo}

\subsection{El modelo}

Se especifica un modelo log-lineal de Poisson \citep{nelderwedderburn1972} para el xG que el equipo $i$ genera contra el equipo $j$:
\begin{equation}
\xG_{ij} \sim \text{Poisson}(\lambda_{ij}), \qquad
\lambda_{ij} \;=\; \mu\, \alpha_i\, \delta_j\, t_{ij},
\label{eq:modelo}
\end{equation}
donde $\alpha_i > 0$ es la fuerza de \emph{generación} de xG del equipo $i$, $\delta_j > 0$ la fuerza de \emph{supresión} de xG del equipo $j$ (valores menores indican mejor defensa), $\mu$ la tasa media del torneo y $t_{ij}$ la exposición del partido en unidades de 90 minutos ($t = 1$ en tiempo reglamentario; $t = 4/3$ con prórroga). No se incluye término de localía: la totalidad del torneo, y en particular la final, se disputa en sedes neutrales.

La ventaja de modelar el xG en lugar del gol es directa. Un modelo de goles requeriría descomponer la tasa en volumen de tiro por probabilidad de conversión, y estimar esta última sobre muestras diminutas ---7 partidos--- introduce una varianza considerable. El xG ya incorpora la calidad de la ocasión: es, por construcción, el número esperado de goles dado el conjunto de disparos generados. El paso de conversión desaparece del modelo.

El modelo es invariante ante la reparametrización $\alpha_i \mapsto c\,\alpha_i$, $\delta_j \mapsto \delta_j/c$. Se impone la restricción de identificabilidad $\sum_i \log\alpha_i = 0$, es decir, media geométrica unitaria de las fuerzas de ataque.

\subsection{Suficiencia de los totales marginales}

En forma logarítmica, el modelo \eqref{eq:modelo} es
\begin{equation}
\log\lambda_{ij} \;=\; \log\mu + \log\alpha_i + \log\delta_j + \log t_{ij},
\end{equation}
un modelo log-lineal con dos factores. Para la familia de Poisson, los estadísticos suficientes de este modelo son los \textbf{totales marginales}: el xG total generado por cada equipo, $\sum_j \xG_{ij}$, y el xG total concedido por cada equipo, $\sum_i \xG_{ij}$. Este resultado es clásico en la teoría de tablas de contingencia \citep{birch1963, demingstephan1940}.

La consecuencia práctica es determinante: \textbf{el ajuste por rival no requiere el xG partido a partido de cada equipo}. Basta conocer los totales agregados de xG a favor y en contra de cada selección, junto con el diseño del torneo ---quién se enfrentó a quién. El modelo utiliza los 102 partidos disputados como estructura de emparejamiento, pero sus parámetros dependen únicamente de los agregados por selección.

La log-verosimilitud, ignorando términos que no dependen de los parámetros, es
\begin{equation}
\ell(\boldsymbol{A}, \boldsymbol{D}) \;=\; \sum_i S_i A_i + \sum_j S^a_j D_j - \sum_{(i,j)\in\mathcal{M}} \mu\, t_{ij}\left(e^{A_i + D_j} + e^{A_j + D_i}\right),
\label{eq:loglik}
\end{equation}
donde $A_i = \log\alpha_i$, $D_j = \log\delta_j$, $S_i$ y $S^a_j$ son los totales marginales de xG observados, y $\mathcal{M}$ el conjunto de los 102 partidos disputados. Cada partido contribuye en ambos sentidos. La maximización se efectuó con L-BFGS-B y gradiente analítico.

Dado que el xG es una magnitud continua no negativa y no un recuento entero, \eqref{eq:loglik} se interpreta como una cuasi-verosimilitud de Poisson: las ecuaciones de estimación dependen solo de la media, de modo que las estimaciones puntuales de $\alpha$ y $\delta$ son válidas aunque el xG no sea estrictamente un conteo.

\subsection{Validación del ajuste}

En el máximo de \eqref{eq:loglik}, las ecuaciones de verosimilitud imponen que los totales marginales ajustados reproduzcan exactamente los observados. Esta identidad constituye un criterio de validación de cumplimiento obligatorio: si los marginales ajustados no coinciden con los observados a precisión de máquina, el optimizador no ha alcanzado el óptimo y las estimaciones no son fiables. La implementación incorpora esta comprobación como aserción que interrumpe la ejecución ante cualquier fallo. En la estimación definitiva, la discrepancia máxima entre marginales ajustados y observados fue de $1{,}4\times 10^{-7}$.

\subsection{Agregación y resolución del partido}
\label{sec:agregacion}

Los goles de penal se añaden como un término aditivo constante, estimado a partir de la frecuencia y la tasa de conversión del torneo:
\begin{equation}
\lambda_{\text{pen}} \;=\; \frac{\text{penales concedidos}}{\text{exposición total}} \times \frac{\text{penales convertidos}}{\text{penales concedidos}} \;=\; \frac{21}{209} \times \frac{15}{21} \;=\; 0{,}072.
\end{equation}

La tasa de goles de cada finalista para la final es
\begin{equation}
\lambda_i \;=\; \mu\, \alpha_i\, \delta_j \;+\; \lambda_{\text{pen}}.
\end{equation}
Los goles de cada equipo se modelan como variables de Poisson independientes, de modo que la matriz de marcadores es el producto exterior de las dos distribuciones marginales.

Al tratarse de una final, un empate en el tiempo reglamentario conduce a prórroga y, de persistir, a lanzamientos desde el punto penal. La prórroga se modela como una continuación con exposición $1/3$ (30 minutos) y tasas $\lambda_i/3$. La tanda de penales se resuelve con probabilidad $0{,}5$ para cada equipo: los datos disponibles no contienen información sobre la habilidad específica en tandas, y asignar ventaja a partir de reputaciones individuales sería opinión, no inferencia.

\section{Resultados}

\subsection{Fuerzas estimadas y efecto del ajuste por rival}

El máximo de \eqref{eq:loglik} se alcanzó con norma del gradiente y discrepancia de marginales del orden de $10^{-7}$. El Cuadro~\ref{tab:fuerzas} presenta las fuerzas de los finalistas antes y después del ajuste por rival.

\begin{table}[htbp]
\centering
\caption{Fuerzas de los finalistas: tasas marginales crudas de xG vs.\ estimaciones ajustadas por rival. Para $\delta$, valores menores indican mejor defensa. Escala: $1{,}00$ = media del torneo.}
\label{tab:fuerzas}
\begin{tabular}{lrrrrrr}
\toprule
& \multicolumn{3}{c}{$\alpha$ (generación de xG)} & \multicolumn{3}{c}{$\delta$ (supresión de xG)} \\
\cmidrule(lr){2-4}\cmidrule(lr){5-7}
& Crudo & Ajustado & Cambio & Crudo & Ajustado & Cambio \\
\midrule
España    & 1{,}339 & \textbf{1{,}844} & $+37{,}7\%$ & 0{,}212 & \textbf{0{,}158} & $-25{,}5\%$ \\
Argentina & 1{,}324 & \textbf{1{,}324} & $\phantom{0}+0{,}0\%$ & 0{,}356 & \textbf{0{,}262} & $-26{,}4\%$ \\
\bottomrule
\end{tabular}
\end{table}

El ajuste modifica las fuerzas de forma reveladora. En ataque, España mejora un 37{,}7\% al corregir por la calidad de sus rivales, mientras que Argentina apenas se mueve: sus adversarios estuvieron muy cerca de la media del torneo. En defensa, ambas mejoran de forma similar. El Cuadro~\ref{tab:caminos} explica el mecanismo.

\begin{table}[htbp]
\centering
\caption{Trayectorias de los finalistas en el torneo.}
\label{tab:caminos}
\begin{tabular}{lll}
\toprule
\textbf{Fase} & \textbf{Rival de España} & \textbf{Rival de Argentina} \\
\midrule
Grupos & Cabo Verde (0--0) & Argelia (3--0) \\
Grupos & Arabia Saudí (4--0) & Austria (2--0) \\
Grupos & Uruguay (1--0) & Jordania (3--1) \\
Dieciseisavos & Austria (3--0) & Cabo Verde (3--2, pr.) \\
Octavos & Portugal (1--0) & Egipto (3--2) \\
Cuartos & Bélgica (2--1) & Suiza (3--1, pr.) \\
Semifinal & Francia (2--0) & Inglaterra (2--1) \\
\bottomrule
\end{tabular}
\end{table}

El cuadro de España incluyó a Uruguay, Portugal, Bélgica y Francia; el de Argentina, a Argelia, Jordania, Cabo Verde, Egipto y Suiza. El ajuste cuantifica esa asimetría sin recurrir a la reputación de los nombres.

\subsection{Ubicación de los finalistas en el torneo}

El Cuadro~\ref{tab:ranking} sitúa a los finalistas frente al resto de selecciones. El resultado matiza el relato dominante en una dirección concreta.

\begin{table}[htbp]
\centering
\caption{Selecciones con mayor fuerza ofensiva ajustada (izquierda) y mejor supresión defensiva ajustada (derecha), por xG. Escala: $1{,}00$ = media del torneo.}
\label{tab:ranking}
\begin{tabular}{lr@{\hskip 1.5cm}lr}
\toprule
\multicolumn{2}{c}{\textbf{Ataque} ($\alpha$, mayor es mejor)} & \multicolumn{2}{c}{\textbf{Defensa} ($\delta$, menor es mejor)} \\
\cmidrule(lr){1-2}\cmidrule(lr){3-4}
& Ajustado & & Ajustado \\
\midrule
Alemania    & 2{,}315 & \textbf{España}    & \textbf{0{,}158} \\
Senegal     & 2{,}170 & \textbf{Argentina} & \textbf{0{,}262} \\
Brasil      & 1{,}957 & Francia            & 0{,}402 \\
Colombia    & 1{,}899 & Canadá             & 0{,}451 \\
Suiza       & 1{,}868 & Uruguay            & 0{,}478 \\
\textbf{España} & \textbf{1{,}844} & México       & 0{,}510 \\
Francia     & 1{,}814 & Brasil             & 0{,}515 \\
Argentina (7\textsuperscript{o}) & 1{,}324 & Congo RD & 0{,}524 \\
\bottomrule
\end{tabular}
\end{table}

Por xG, \textbf{ni España ni Argentina son el mejor ataque del torneo}: España figura sexta y Argentina bastante más abajo. El ataque más potente, medido por generación de xG ajustada, es el de Alemania, eliminada en dieciseisavos por Paraguay en la tanda de penales. En cambio, \textbf{en defensa los finalistas ocupan las dos primeras posiciones}: España concede el menor xG del torneo y Argentina el segundo menor. La final enfrenta a las dos mejores defensas del campeonato, y la ventaja de España es esencialmente defensiva.

\subsection{Predicción}

Con $\mu = 1{,}434$, el xG esperado para la final es
\begin{align}
\E[\xG_{\text{ESP}}] &= \mu\, \alpha_{\text{ESP}}\, \delta_{\text{ARG}} = 1{,}434 \times 1{,}844 \times 0{,}262 = 0{,}69, \\
\E[\xG_{\text{ARG}}] &= \mu\, \alpha_{\text{ARG}}\, \delta_{\text{ESP}} = 1{,}434 \times 1{,}324 \times 0{,}158 = 0{,}30,
\end{align}
y sumando el término de penales,
\begin{align}
\lambda_{\text{ESP}} &= 0{,}69 + 0{,}072 = \mathbf{0{,}765}, \\
\lambda_{\text{ARG}} &= 0{,}30 + 0{,}072 = \mathbf{0{,}372}.
\end{align}

El total esperado, $1{,}14$ goles, se sitúa muy por debajo de la media del torneo ($2{,}73$): el modelo anticipa una final extremadamente cerrada, consecuencia directa de que ambos finalistas son los mejores supresores de xG del campeonato.

\begin{table}[htbp]
\centering
\caption{Probabilidades de resultado.}
\label{tab:pred}
\begin{tabular}{lrr}
\toprule
& \textbf{90 minutos} & \textbf{Campeón} \\
\midrule
España    & 41{,}7\% & \textbf{64{,}9\%} \\
Empate    & 41{,}9\% & --- \\
Argentina & 16{,}5\% & \textbf{35{,}1\%} \\
\bottomrule
\end{tabular}
\end{table}

La elevada probabilidad de empate en el tiempo reglamentario (41{,}9\%) es consecuencia de las tasas bajas: cuando ambas medias son pequeñas, la masa de probabilidad se concentra en los marcadores mínimos, muchos de ellos igualados. El modelo asigna en consecuencia una probabilidad sustancial (41{,}9\%) de que la final se prolongue a la prórroga.

\subsection{Matriz de marcadores}

\begin{table}[htbp]
\centering
\caption{Distribución de marcadores en el tiempo reglamentario (\%). Filas: goles de España; columnas: goles de Argentina. La diagonal corresponde a los empates.}
\label{tab:matriz}
\begin{tabular}{c|rrrrr}
\toprule
\diagbox[width=2.2cm]{ESP}{ARG} & \textbf{0} & \textbf{1} & \textbf{2} & \textbf{3} & \textbf{4} \\
\midrule
\textbf{0} & \cellcolor{COLNEU!35}\textbf{32{,}07} & \cellcolor{COLARG!25}11{,}93 & \cellcolor{COLARG!10}2{,}22 & 0{,}28 & 0{,}03 \\
\textbf{1} & \cellcolor{COLESP!40}24{,}54 & \cellcolor{COLNEU!22}9{,}13 & \cellcolor{COLARG!8}1{,}70 & 0{,}21 & 0{,}02 \\
\textbf{2} & \cellcolor{COLESP!20}9{,}39 & \cellcolor{COLESP!10}3{,}49 & \cellcolor{COLNEU!6}0{,}65 & 0{,}08 & 0{,}01 \\
\textbf{3} & \cellcolor{COLESP!7}2{,}40 & \cellcolor{COLESP!5}0{,}89 & 0{,}17 & 0{,}02 & 0{,}00 \\
\textbf{4} & 0{,}46 & 0{,}17 & 0{,}03 & 0{,}00 & 0{,}00 \\
\bottomrule
\end{tabular}
\end{table}

El marcador más probable es el empate sin goles (32{,}07\%), seguido de España 1--0 (24{,}54\%). El Cuadro~\ref{tab:mercados} recoge algunos agregados derivados.

\begin{table}[htbp]
\centering
\caption{Probabilidades derivadas de la matriz de marcadores.}
\label{tab:mercados}
\begin{tabular}{lr}
\toprule
\textbf{Suceso} & \textbf{Probabilidad} \\
\midrule
Menos de 2{,}5 goles & 89{,}3\% \\
Más de 2{,}5 goles & 10{,}7\% \\
Ambos equipos marcan & 16{,}6\% \\
España mantiene la portería a cero & 68{,}9\% \\
Argentina mantiene la portería a cero & 46{,}5\% \\
El partido llega a la prórroga & 41{,}9\% \\
\bottomrule
\end{tabular}
\end{table}

\subsection{Cuantificación de la incertidumbre}

Se efectuó un \textit{bootstrap} paramétrico de 4.000 réplicas, remuestreando el xG total generado y concedido por cada finalista según su distribución de muestreo y recalculando la predicción completa en cada réplica.

\begin{table}[htbp]
\centering
\caption{\textit{Bootstrap} paramétrico, 4.000 réplicas.}
\label{tab:boot}
\begin{tabular}{lrr}
\toprule
& \textbf{Mediana} & \textbf{IC 95\%} \\
\midrule
P(España campeón) & 63{,}7\% & [35{,}6\%;\ 89{,}1\%] \\
\bottomrule
\end{tabular}
\end{table}

España resulta favorita en el 83{,}2\% de las réplicas. El intervalo es ancho: con 7 partidos por selección no cabe mayor precisión, y el uso del xG ---más estable que el gol--- reduce el ruido pero no lo elimina. El límite inferior del intervalo alcanza el 35{,}6\%, de modo que el modelo no descarta que Argentina sea, en realidad, ligeramente favorita.

\section{Interpretación}

\subsection{Por qué el relato dominante no predice}

La caracterización ``el mejor ataque contra la mejor defensa'' es correcta como descripción: Argentina anotó más goles que nadie y España concedió menos que nadie. Pero el xG revela que ambos registros descansan en gran medida sobre la varianza de finalización. En generación de xG, ambos finalistas son casi idénticos ($1{,}92$ por 90 minutos), y ninguno de los dos es el ataque más peligroso del torneo. La diferencia entre 18 y 12 goles no proviene de crear mejores ocasiones ---las crean equivalentes--- sino de convertirlas por encima o por debajo del valor esperado, un componente que no es predictivo.

\subsection{Dónde reside realmente la ventaja de España}

El xG reubica la ventaja de España en su lugar correcto: la defensa. España concede el menor xG por 90 minutos de todo el torneo ($0{,}305$), con Argentina en segundo lugar ($0{,}512$). Esta superioridad es genuina y estable, a diferencia del ``un gol encajado'', que combinaba la excelencia defensiva con una dosis de fortuna. La final, por xG, es un duelo entre las dos mejores defensas del campeonato, en el que España parte con una ventaja defensiva medible.

\subsection{Sobre la magnitud del resultado}

Una probabilidad de 64{,}9\% no es la predicción de que España ganará. Significa que, si esta final se disputara 100 veces bajo las condiciones que el modelo captura, Argentina levantaría el trofeo en aproximadamente 35 ocasiones. El límite inferior del intervalo de confianza alcanza el 35{,}6\%: el modelo no descarta que ambas selecciones estén sustancialmente igualadas, o incluso que Argentina sea ligeramente favorita. El partido se disputa una sola vez.

\section{Limitaciones}

\begin{enumerate}
\item \textbf{El xG procede de un modelo propietario.} El xG de Sofascore descansa en una especificación no pública: se desconoce qué variables emplea y cómo fue entrenado. Es una medición valiosa producida por especialistas con acceso a datos posicionales masivos, pero no una verdad de campo, sino la salida del modelo de un tercero. Un modelo de xG propio, entrenado y auditado por el analista sobre los datos de disparo con coordenadas obtenidos aquí (2.591 disparos), permitiría un control metodológico completo, si bien excede el alcance temporal de este trabajo.

\item \textbf{Ausencia de la corrección de \citet{dixoncoles1997} para marcadores bajos.} El término $\tau$ corrige la dependencia empíricamente observada en los resultados $0$--$0$, $1$--$0$, $0$--$1$ y $1$--$1$, que el Poisson bivariado independiente no captura. Su estimación exige datos partido a partido. Con tasas tan bajas como las aquí estimadas, es probable que el modelo subestime aún más la probabilidad de empate y, por tanto, de prórroga. Un Poisson bivariado con covarianza explícita \citep{karlisntzoufras2003} abordaría la misma limitación.

\item \textbf{Sin información de plantilla.} El modelo desconoce lesiones, sanciones y estado de forma individual. En particular, no incorpora que España afronta la final sin algunos de sus jugadores clave en su mejor momento.

\item \textbf{Tanda de penales con probabilidad $0{,}5$.} Los datos no contienen información sobre habilidad en tandas.

\item \textbf{Estacionariedad supuesta.} Las fuerzas se estiman como constantes a lo largo del torneo. \citet{dixoncoles1997} proponen una ponderación exponencial decreciente en el tiempo, no aplicada aquí por la brevedad de la ventana (cinco semanas).
\end{enumerate}

\section{Conclusiones}

\begin{enumerate}
\item \textbf{El relato ``mejor ataque contra mejor defensa'' describe el pasado pero no predice el futuro.} Por xG, ambos finalistas generan peligro de forma casi idéntica ($1{,}924$ frente a $1{,}920$ de xG por 90 minutos) y ninguno es el ataque más peligroso del torneo. La disparidad entre 18 y 12 goles es varianza de finalización, no diferencia en la creación de ocasiones.

\item \textbf{El xG resuelve la paradoja del gol único de España sin correcciones \textit{ad hoc}.} Donde el recuento de goles ofrecía un estimador degenerado ($0{,}14$ goles concedidos por partido), el xG ofrece un valor estable y bajo ($0{,}305$ por 90 minutos): una defensa excelente, no sobrenatural.

\item \textbf{La ventaja de España es defensiva.} España concede el menor xG por 90 minutos del torneo y Argentina el segundo menor. La final enfrenta a las dos mejores defensas del campeonato, con España en ventaja medible.

\item \textbf{El ajuste por rival favorece a España, sobre todo en ataque.} Su fuerza ofensiva mejora un 37{,}7\% al corregir por la calidad de sus adversarios, mientras que la de Argentina no se mueve. La trayectoria de España ---Uruguay, Portugal, Bélgica, Francia--- fue sustancialmente más exigente.

\item \textbf{La suficiencia de los totales marginales permitió el ajuste por rival a partir de agregados por selección.} El resultado clásico sobre modelos log-lineales de Poisson \citep{birch1963} hizo depender la estimación únicamente de los totales de xG a favor y en contra, más el diseño del torneo.

\item \textbf{La predicción asigna a España un 64{,}9\% de probabilidad de título} (IC 95\%: [35{,}6\%; 89{,}1\%]), con $\lambda_{\text{ESP}} = 0{,}77$ y $\lambda_{\text{ARG}} = 0{,}37$, y anticipa una final muy cerrada (1{,}14 goles esperados, con un 41{,}9\% de probabilidad de prórroga). El marcador más probable es el empate sin goles (32{,}1\%).
\end{enumerate}

\section*{Apéndice: reproducibilidad}
\addcontentsline{toc}{section}{Apéndice}

\textbf{Datos.} xG de Sofascore (competición \texttt{FIFA World Cup}, temporada \texttt{2026}), recuento de goles y calendario de FBref. Extracción: 17--18 de julio de 2026, tras 102 de 104 partidos.

\textbf{Entorno.} Python 3.12; \texttt{ScraperFC} 4.5.0, \texttt{pandas}, \texttt{numpy}, \texttt{scipy}.

\textbf{Archivos.}
\begin{itemize}
\item \texttt{sofascore\_xg.py} --- extracción de los mapas de disparo y agregación de xG por partido.
\item \texttt{modelo\_xg.py} --- estimación de fuerzas ajustadas por rival, con la aserción de validación sobre los marginales, y predicción.
\item \texttt{master\_xg.csv}, \texttt{calendario.csv}, \texttt{fuerzas\_xg.csv} --- datos intermedios y fuerzas estimadas de las 48 selecciones.
\end{itemize}

\textbf{Nota sobre el xG y las tandas de penales.} El campo de resultado de Sofascore incluye los goles de las tandas de penales. El recuento de goles se tomó de FBref para evitar esta contaminación; el xG, no afectado por las tandas, se tomó de Sofascore.

\bibliography{referencias}

\end{document}
